\chapter{Preliminaries}
\section{Categories as deductive systems}
Following \cite{LambekScott1994}, we will define a category as a particular kind of graph, but first we'll construct a series of structures that will provide us
with a working intuitionistic deductive system made entirely of objects and arrows.

\begin{definition}[Graph]
    A \textit{graph} $\mathcal G$ is given by a class of \textit{objects}, $\ob(\mathcal G)$, a class of \textit{arrows}, $\mor(\mathcal G)$, and two functions
    \textit{source} and \textit{target}
    \[ s,t:\mor(\mathcal G)\to\ob(\mathcal G),\]
    which assign to each arrow $f$ its source $s(f)$ and its target $t(f)$.
\end{definition}

\begin{notation}
    We will denote by $\Hom_{\mathcal G}(A,B)$, omitting the subscript when it is clear from the context, the class of arrows $f$ such that $s(f)=A$ and $t(f)=B$.
\end{notation}

\begin{definition}[Deductive system]
    A \textit{deductive system} $\mathcal D$ is a graph in which, for each object $A$, there is an arrow $1_A : A \to A$, called \textit{identity arrow}, and there exists a
    \textit{composition} function between arrows
    \[ \circ : (f,g) \mapsto g \circ f \qquad \text{for all $f$ and $g$ such that $t(f)=s(g)$.} \]
    We could also define the composition as a family of functions, one for each triple of objects $A,B,C \in \ob(\mathcal D)$:
    \[ \circ_{A,B,C} : \Hom(A,B) \times \Hom(B,C) \to \Hom(A,C) : (f,g) \mapsto g \circ f, \]
    and we would get an equivalent definition.
\end{definition}

\begin{notation}
    For brevity, we will often write the composition $g\circ f$ of two arrows $f$ and $g$ as $gf$. Following a familiar logical notation, we will also denote the composition as:
    \[ \frac{A \xrightarrow{f} B \quad B \xrightarrow{g} C}{A \xrightarrow{gf} C}. \]
\end{notation}

The objects of a deductive system can be thought of as \textit{propositions}, the arrows as \textit{proofs} and the composition as a \textit{rule of inference}.

\begin{definition}[Conjunction calculus]
    A \textit{conjunction calculus} is a deductive system $\mathcal C$, with a special object $\top$, the \textit{tautological proposition}, and a function
    \[ \wedge : \ob(\mathcal C)\times\ob(\mathcal C) \to \ob(\mathcal C) : (A,B) \mapsto A\wedge B, \]
    called \textit{conjunction}. We also have, in addition, that the following rules, which state the existence of certain classes of arrows, hold:
    \begin{itemize}
        \item[(\textnormal{Conj 1})] for each object $A$ there exists an arrow $!_A : A \to \top$;
        \item[(\textnormal{Conj 2a})] for each pair of objects $A,B$ there exists an arrow $\pi_{A,B} : A \wedge B \to A$;
        \item[(\textnormal{Conj 2b})] for each pair of objects $A,B$ there exists an arrow $\pi_{A,B}' : A \wedge B \to B$;
        \item[(\textnormal{Conj 2c})] for each pair of arrows $f : A \to B$ and $g : A \to C$ there exists an arrow $\langle f,g\rangle : A \to B \wedge C$, the \textit{pairing arrow}, which can also be written as:
        \[ \frac{A \xrightarrow{f} B \quad A \xrightarrow{g} C}{A \xrightarrow{\langle f,g\rangle} B \wedge C}.
        \]
    \end{itemize}
\end{definition}

Finding a proof in a conjunction calculus means finding an arrow that starts from the source proposition and ends at the target proposition.

\begin{example}[Commutativity of conjunction]
    Given two objects $A$ and $B$ in a conjunction calculus, we can find a proof of $B \land A$ from $A \land B$ using the projections and the pairing arrow
    \[ \frac{A \wedge B \xrightarrow{\pi_{A,B}'} B \quad A \wedge B \xrightarrow{\pi_{A,B}} A}{A \wedge B \xrightarrow{\langle \pi_{A,B}',\pi_{A,B}\rangle} B \wedge A}, \]
    therefore the arrow $\langle \pi_{A,B}',\pi_{A,B}\rangle : A \wedge B \to B \wedge A$ is a proof of $B \land A$ from $A \land B$.
\end{example}

\begin{example}[Associativity of conjunction]\label{ex:associativity-of-conjunction}
    Given three objects $A$, $B$ and $C$ in a conjunction calculus, we can a proof of $A \land (B \land C)$ from $(A \land B) \land C$ is:
    \[ \alpha_{A,B,C} := \langle \pi_{A,B}\circ \pi_{A \land B,C}, \langle \pi_{A,B}'\circ \pi_{A \land B,C}, \pi_{A \land B,C}'\rangle\rangle. \]
\end{example}

Continuing this analogy, one could think of (Conj 2a) and (Conj 2b) as the \textit{elimination rules} for conjunction in the natural deduction system \cite{ChiswellHodgesAppendixA,Mennuni2025}, of (Conj 2c) as the \textit{introduction rule} for conjunction, and of
(Conj 1) as the \textit{introduction rule} for $\top$ united to \textit{weakening} rule. We will progressively add more rules, or rather more arrows, to our deductive systems until we reach a complete intuitionistic deductive system.

\begin{remark}
    Using the rules of a conjunction calculus, we can also define new rules of inference, or rather new operations on arrows, using the existing ones (for now only the composition and the pairing).
    For instance, given two arrows $f : A \to B$ and $g : C \to D$, we can define the arrow $f \wedge g : A \wedge C \to B \wedge D$ by:
    \[ \frac{\frac{A \wedge C \xrightarrow{\pi_{A,C}} A \quad A \xrightarrow{f} B}{A \wedge C \xrightarrow{f\pi_{A,C}} B} \quad \frac{A \wedge C \xrightarrow{\pi_{A,C}'} C \quad C \xrightarrow{g} D}{A \wedge C \xrightarrow{g\pi_{A,C}'} D}}{A \wedge C \xrightarrow{\langle f\pi_{A,C},g\pi_{A,C}'\rangle} B \wedge D}, \]
    and then we put $f \wedge g := \langle f\pi_{A,C},g\pi_{A,C}'\rangle$. Thus, from now on, we will simply write:
    \[ \frac{A \xrightarrow{f} B \quad C \xrightarrow{g} D}{A \wedge C \xrightarrow{f \wedge g} B \wedge D}. \]
\end{remark}

We're still far from having a complete intuitionistic deductive system, we still don't have a way to express the implication for example, so we give the following definition.

\begin{definition}[Positive intuitionistic calculus]
    A \textit{positive intuitionistic propositional calculus} is a conjunction calculus $\mathcal P$ with a function
    \[ \Leftarrow : \ob(\mathcal P)\times\ob(\mathcal P) \to \ob(\mathcal P) : (A,B) \mapsto A \Leftarrow B, \]
    called \textit{implication}, and the following rules of inference:
    \begin{itemize}
        \item[(\textnormal{Impl 1})] for each pair of objects $A$ and $B$ there exists an arrow $\ev_{A,B} : (A \Leftarrow B) \wedge B \to A$, called \textit{evaluation arrow};
        \item[(\textnormal{Impl 2})] for each triple of objects $A$, $B$ and $C$ and for each arrow $h : A \wedge B \to C$ there is an arrow $h^* : A \to (C \Leftarrow B)$, called \textit{currying of $h$}.
    \end{itemize}
\end{definition}

\begin{remark}
    As for the conjunction rules, we can also see here the analogy with the natural deduction system, in fact (Impl 1) is the \textit{elimination rule} for implication, however (Impl 2) does not correspond to any basic inference rule, but rather to a derived one:
    \[ \textnormal{In}_{\Leftarrow}\frac{\displaystyle\textnormal{El}_{\Leftarrow}\frac{\displaystyle \textnormal{Wk}\frac{\displaystyle \textnormal{In}_{\Leftarrow}\frac{\displaystyle\frac{}{A \land B \vdash C}}{\vdash C \Leftarrow A \land B}}{A,B \vdash C \Leftarrow A \land B} \quad \displaystyle\textnormal{In}_{\land}\frac{\displaystyle\textnormal{Wk}\frac{\displaystyle\textnormal{Ax}\frac{}{A \vdash A}}{A, B \vdash A} \quad \displaystyle\textnormal{Wk}\frac{\displaystyle\textnormal{Ax}\frac{}{B \vdash B}}{A, B \vdash B}}{A, B \vdash A \land B}}{A,B \vdash C}}{A \vdash C \Leftarrow B}. \]
\end{remark}

\begin{remark}[Derived rules of inference]
    From (Impl 1) and (Impl 2) we can derive the following rules of inference:
    \begin{itemize}
        \item[(\textnormal{Impl 1'})] for each pair of objects $B$ and $C$ there exists an arrow $\eta_{C,B} : C \to (C \land B) \Leftarrow B$, called \textit{co-evaluation arrow};
        \item[(\textnormal{Impl 2'})] given an arrow $g : D \to A$ and an object $B$, we have the \textit{co-currying} arrow
        \[ \frac{D \xrightarrow{g} A}{(D \Leftarrow B) \xrightarrow{g \Leftarrow 1_B} (A \Leftarrow B)}, \]
    which expresses the covariance of implication in its first argument.   
    \end{itemize}
    Both of these alternative rules can be derived from the original ones, for example, by taking $\eta_{C,B} := 1_{C \land B}^*$ and $g \Leftarrow 1_B := (g \circ \ev_{D,B})^*$. It can also be shown that (Impl 1') and (Impl 2') imply (Impl 2),
    but not (Impl 1), so they are not equivalent to the original rules.
\end{remark}

\begin{notation}
    We will often forget the subscripts of the evaluation and co-evaluation arrows when they are clear from the context, denoting them simply by $\ev$ and $\eta$, respectively.
\end{notation}

Before moving on, let's also define two more rules of inference which follow from the previous ones and will be very useful in the next sections.

\begin{example}
    Given an arrow $f : A \to B$ we can define the arrow $\ulcorner f \urcorner$ as
    \[ \frac{A \xrightarrow{f} B}{\top \xrightarrow{\ulcorner f \urcorner} B \Leftarrow A}, \]
    that is just $\ulcorner f \urcorner := (f \circ \pi_{\top,A}')^*$. Similarly, given $g : \top \to B \Leftarrow A$, we can define 
    \[ \frac{\top \xrightarrow{g} B \Leftarrow A}{A \xrightarrow{g^\sharp} B} \]
    with $g^\sharp := \ev_{B,A} \circ \langle g\circ !_A,1_A\rangle$. 
\end{example}

We now need to add falsehood and disjunction to our positive intuitionistic propositional calculus, to obtain a complete intuitionistic propositional calculus.
To this end we give the following definition.

\begin{definition}[Intuitionistic propositional calculus]
    An \textit{intuitionistic propositional calculus} is a positive intuitionistic propositional calculus $\mathcal I$ with a special object $\bot$, the \textit{falsehood formula}, a function
    \[ \lor : \ob(\mathcal I)\times\ob(\mathcal I) \to \ob(\mathcal I) : (A,B) \mapsto A \lor B, \]
    called \textit{disjunction}, and the following rules of inference:
    \begin{itemize}
        \item[(\textnormal{Int 1})] for each object $A$ there exists an arrow $?_A : \bot \to A$;
        \item[(\textnormal{Int 2a})] for each pair of objects $A$ and $B$ there exists an arrow $\iota_{A,B} : A \to A \lor B$;
        \item[(\textnormal{Int 2b})] for each pair of objects $A$ and $B$ there exists an arrow $\iota_{A,B}' : B \to A \lor B$;
        \item[(\textnormal{Int 2c})] for each triple of objects $A$, $B$ and $C$ there exists an arrow
        \[ \zeta_{A,B}^C : (C \Leftarrow A) \land (C \Leftarrow B) \to C \Leftarrow (A \lor B),\]
        called \textit{disjunction elimination arrow}.
    \end{itemize}
\end{definition}

\begin{remark}
    One can recognize in (Int 2a) and (Int 2b) the introduction rules for disjunction in the natural deduction system,
    in (Int 2c) an internalized version of the elimination rule for disjunction, and in (Int 1) the elimination rule for falsehood.
\end{remark}

Let's define another useful inference rule that can be derived in this system.

\begin{example}[Co-pairing arrow]
    Given the arrows $f : A \to C$ and $g : B \to C$, we can define the \textit{co-pairing arrow}
    \[ \frac{A \xrightarrow{f} C \quad B \xrightarrow{g} C}{A \lor B \xrightarrow{[f,g]} C}, \]
    where $[f,g] := (\zeta_{A,B}^C \circ \langle \ulcorner f \urcorner, \ulcorner g \urcorner\rangle)^\sharp$.
    This inference rule is the direct counterpart of the elimination rule for disjunction in $\mathsf{ND}$ (the Natural Deduction system). 
\end{example}

We have now constructed a complete intuitionistic propositional calculus, to conclude this section we will take the final step and define a category that has the structure of classical propositional logic.

\begin{definition}[Classical propositional calculus]
    A \textit{classical propositional calculus} is an intuitionistic propositional calculus $\mathcal C$ where for each object $A$ there exists an arrow:
    \[ \neg\neg : \bot \Leftarrow (\bot \Leftarrow A) \to A. \]
\end{definition}

Although intuitionistic and classical propositional calculi will not play an important role in the rest of this work, we have defined them to show how the structure of a category
can be used to represent a complete deductive system for classical and intuitionistic propositional logic, thinking of the objects of the category as propositions and of the arrows as proofs.

\section{Cartesian closed categories}
Now that we have defined enough structure to represent a complete deductive system, we want to impose some rules on the arrows of a category to obtain a structure that is more suitable for our purposes, that is a cartesian closed structure.

\begin{definition}[Category]
    A \textit{category} is a deductive system $\mathcal C$ where the following equations hold:
    \[ f \circ 1_A = f, \qquad 1_B \circ f = f, \qquad (h \circ g) \circ f = h \circ (g \circ f),\]
    for all arrows $f : A \to B$, $g : B \to C$ and $h : C \to D$. The first two equations are called \textit{identity laws}, while the third one is the \textit{associativity law}.
\end{definition}

Let's observe that we can always construct a category from a deductive system, this notion will be developed more in depth later.

\begin{remark}[Category freely generated by a deductive system]\label{rem:category-generated-by-deductive-system}
    Given a deductive system $\mathcal D$ we can always construct the \textit{freely generated category} $\mathcal C$ over $\mathcal D$, by taking the same objects and, as arrows, the equivalence classes of arrows of $\mathcal D$ modulo the identity and associativity laws.
    In particular, we can define
    \[ \ob(\mathcal C) = \ob(\mathcal D), \quad \mor(\mathcal C) = \mor(\mathcal D)/\sim_{\textnormal{Cat}}, \]
    where $\sim_{\textnormal{Cat}}$ is the equivalence relation generated by the identity and associativity laws which is compatible with composition,
    i.e.,\ the smallest equivalence relation on parallel arrows such that the equations above hold (it is well-defined, for example, since there exists the equivalence relation of having the same source and target, which is compatible with composition and contains the identity and associativity laws). 
\end{remark}

A cartesian structure over a category basically means that the category resembles in some sense the category of sets, in particular it has a terminal object and it is closed under binary products,
we'll use the following definition, which is equivalent to the usual one, but algebraic in nature, in the sense that it is based on the existence of certain equations that the arrows of the category must satisfy.

\begin{definition}[Cartesian category]
    A \textit{cartesian category} is a category $\mathcal C$ that is also a conjunction calculus, where the following equations hold:
    \begin{itemize}
        \item[(\textnormal{CC1})] for every object $A$ and every arrow $f : A \to \top$, we have $f = !_A$;
    \end{itemize}
    for $f : C \to A$, $g : C \to B$ and $h : C \to A \land B$ we have:
    \begin{itemize}
        \item[(\textnormal{CC2a})] $\pi_{A,B} \circ \langle f,g\rangle = f$;
        \item[(\textnormal{CC2b})] $\pi_{A,B}' \circ \langle f,g\rangle = g$;
        \item[(\textnormal{CC2c})] $\langle \pi_{A,B} \circ h, \pi_{A,B}' \circ h\rangle = h$.
    \end{itemize}
\end{definition}

The condition (CC1) states that $\top$ is the terminal object of the category, while (CC2a), (CC2b) and (CC2c) state that $A \land B$ 
satisfies the universal property of the product of $A$ and $B$, with $\pi_{A,B}$ and $\pi_{A,B}'$ as projections.

\begin{notation}[Terminal object and product]
    From now on, we will denote by $1$ the terminal object of a cartesian category, and by $A \times B$ the product of two objects $A$ and $B$.
\end{notation}

\begin{remark}[Equivalent definition of terminal object]
    The equation (CC1) is equivalent to the following
    \begin{itemize}
        \item[(\textnormal{CC1'})] for all $f : A \to B$ we have $!_B \circ f = !_A$ and $1_1 = !_1$.
    \end{itemize}
    Indeed, one can easily show that (CC1) implies (CC1') and vice versa, operating in the metatheory; the core difference is in the fact that (CC1') states the uniqueness algebraically, while (CC1) states it directly.
    There's no difference in terms of choice, since the conjunction calculus already provides us with the collection of arrows $!_A$ for each object $A$.
\end{remark}

As a consequence of the definition we have the following useful lemma.

\begin{lemma}[Distributive law]
    In a cartesian category, for all $f : C \to A$, $g : C \to B$ and $h : D \to C$ we have:
    \[ \langle f,g\rangle \circ h = \langle f \circ h, g \circ h\rangle. \]
\end{lemma}

\begin{proof}
    The proof is just a chain of equalities using the definition of cartesian category:
    \begin{align*}
        \langle f,g\rangle \circ h &= \langle \pi_{A,B} \circ (\langle f,g\rangle \circ h), \pi_{A,B}' \circ (\langle f,g\rangle \circ h)\rangle &&(\textnormal{by CC2c}) \\
                                   &= \langle (\pi_{A,B} \circ \langle f,g\rangle) \circ h, (\pi_{A,B}' \circ \langle f,g\rangle) \circ h\rangle &&(\textnormal{by associativity law}) \\
                                   &= \langle f \circ h, g \circ h\rangle &&(\textnormal{by CC2a and CC2b}).
    \end{align*}
\end{proof}

The distributive law immediately implies the fact that $\cdot \times \cdot$ is a functor.

\begin{notation}[Product of arrows]
    Given $f : A \to B$ and $g : C \to D$, we will write $f \times g : A \times C \to B \times D$ for the arrow $\langle f \circ \pi_{A,C}, g \circ \pi_{A,C}'\rangle$.
\end{notation}

\begin{corollary}[Functoriality of the product of arrows]
    In a cartesian category $\mathcal C$, the operation (inference rule) $\times$ is a functor $\times : \mathcal C \times \mathcal C \to \mathcal C$.
\end{corollary}

\begin{proof}
    We need to check that $\times$ preserves identities and composition.
    \begin{itemize}
        \item[$\boxed{\textnormal{identity}}$] Given $A$ and $B$ in $\mathcal C$, we have:
        \begin{align*}
            1_A \times 1_B &= \langle 1_A \circ \pi_{A,B}, 1_B \circ \pi_{A,B}'\rangle &&(\textnormal{by definition of $\times$}) \\
                           &= \langle \pi_{A,B}, \pi_{A,B}'\rangle &&(\textnormal{by identity law}) \\
                           &= \langle \pi_{A,B} \circ 1_{A \times B}, \pi_{A,B}' \circ 1_{A \times B}\rangle &&(\textnormal{by identity law}) \\
                           &= 1_{A \times B} &&(\textnormal{by CC2c}).
        \end{align*} 
        \item[$\boxed{\textnormal{composition}}$] Given $f : A \to B$, $g : C \to D$, $f' : B \to B'$ and $g' : D \to D'$, similar steps give us:
        \begin{align*}
            (f' \times g') \circ (f \times g) &= \langle f' \circ \pi_{B,D}, g' \circ \pi_{B,D}'\rangle \circ (f \times g) &&(\textnormal{by definition of $\times$}) \\
                                              &= \langle (f' \circ \pi_{B,D}) \circ (f \times g), (g' \circ \pi_{B,D}') \circ (f \times g)\rangle &&(\textnormal{by distributive law})\\
                                              &= \langle (f' \circ f) \circ \pi_{A,C}, (g' \circ g) \circ \pi_{A,C}'\rangle &&(\textnormal{by CC2a and CC2b}) \\
                                              &= (f' \circ f) \times (g' \circ g) &&(\textnormal{by definition of $\times$}).
        \end{align*}
    \end{itemize}
\end{proof}

We are now ready to give the definition of a cartesian closed category, which will be the central structure of this work. The idea is that a cartesian closed category is a category that behaves like $\mathsf{Set}$ in many respects.

\begin{definition}[Cartesian closed category]
    A \textit{cartesian closed category} is a cartesian category $\mathcal C$ that is also a positive intuitionistic propositional calculus, where the following equations hold:
    \begin{itemize}
        \item[(\textnormal{CCC1})] for all $h : A \times B \to C$ we have:
        \[ \ev_{C,B} \circ (h^* \times 1_B) = \ev_{C,B} \circ \langle h^* \circ \pi_{A,B}, \pi_{A,B}'\rangle = h; \]
        \item[(\textnormal{CCC2})] for all $k : A \to C \Leftarrow B$ we have:
        \[ (\ev_{C,B} \circ (k \times 1_B))^* = (\ev_{C,B} \circ \langle k \circ \pi_{A,B}, \pi_{A,B}'\rangle)^* = k. \]
    \end{itemize} 
\end{definition}

Even though the definition just given seems a bit abstract, we can give an equivalent, more concrete, definition of cartesian closed category as follows.

\begin{remark}[Equivalent definition of cartesian closed category]
    If we define the functions:
    \begin{align*}
        & \Phi : \Hom(A \times B, C) \to \Hom(A, C \Leftarrow B) : h \mapsto h^*, \\
        & \Psi : \Hom(A, C \Leftarrow B) \to \Hom(A \times B, C) : k \mapsto \ev_{C,B} \circ (k \times 1_B),
    \end{align*}
    then (CCC1) is equivalent to $\Psi \circ \Phi(h) = h$ for all $h : A \times B \to C$, while (CCC2) is equivalent to $\Phi \circ \Psi(k) = k$ for all $k : A \to C \Leftarrow B$.
\end{remark}

The second definition is exactly the universal property of the exponential object, that is $\Hom(A \times B, C) \cong \Hom(A, C^B)$. Therefore we can also say that a cartesian closed category is a cartesian category that is closed under exponentials.

\begin{notation}[Exponential object]
    From now on, we will denote by $B^A$ the exponential object of $B$ and $A$ in a cartesian closed category.
\end{notation}

\begin{remark}[Adjunction between $- \times B$ and $-^B$]
    The universal property of the exponential object can be equivalently expressed as an adjunction between the functors $- \times B$ and $-^B$. This
    perfectly captures the idea of currying and uncurrying, that is the idea that a function of two variables can be seen as a function of one variable that returns a function of one variable.
\end{remark}

Let's now state two results that will be useful in the next sections.

\begin{proposition}[Distributivity of currying]
    In a cartesian closed category, given $h \in \Hom(A \times B, C)$ and $k \in \Hom(D, A)$, we have:
    \[ h^* \circ k = (h \circ (k \times 1_B))^*. \]
\end{proposition}

\begin{proof}
    We omit composition symbols for brevity:
    \begin{align*}
        h^* k &= (\ev_{C,B} (h^*k \times 1_B))^* &&(\textnormal{by CCC2}) \\
              &= (\ev_{C,B} (h^* \times 1_B) (k \times 1_B))^* &&(\textnormal{by functoriality of $\times$}) \\
              &= (h (k \times 1_B))^* &&(\textnormal{by CCC1}).
    \end{align*}
\end{proof}

\begin{proposition}[Morphisms as global elements]\label{prop:morphisms-as-global-elements}
    In any cartesian closed category, given two objects $A$ and $B$, it holds that:
    \[ \Hom(A,B) \cong \Hom(1, B^A). \]
\end{proposition}

\begin{proof}
    We've already introduced the functions
    \[ \ulcorner \cdot \urcorner : \Hom(A,B) \to \Hom(1, B^A), \qquad (\cdot)^\sharp : \Hom(1, B^A) \to \Hom(A,B),\]
    defined as $\ulcorner f \urcorner := (f \circ \pi_{1,A}')^*$ and $g^\sharp := \ev_{B,A} \circ \langle g\circ !_A,1_A\rangle$, respectively.
    We need to show that they are inverses of each other. First, given $f : A \to B$, we have:
    \begin{align*}
        \ulcorner f \urcorner^\sharp &= \ev_{B,A} \circ \langle (f \circ \pi_{1,A}')^* \circ !_A, 1_A\rangle &&(\textnormal{by the definitions}) \\
                                     &= \ev_{B,A} \circ [((f \circ \pi_{1,A}')^* \times 1_A) \circ \langle !_A, 1_A\rangle] &&(\textnormal{by $\star$}) \\
                                     &= [\ev_{B,A} \circ ((f \circ \pi_{1,A}')^*\times 1_A)] \circ \langle !_A, 1_A\rangle &&(\textnormal{by associativity law}) \\
                                     &= (f \circ \pi_{1,A}') \circ \langle !_A, 1_A\rangle &&(\textnormal{by CCC1}) \\
                                     &= f &&(\textnormal{by associativity law + CC2b}),
    \end{align*}
    where $\star$ denotes the following identity, given $\ell \in \Hom(1,B^A)$:
    \[ \langle \ell \circ !_A, 1_A\rangle = (\ell \times 1_A) \circ \langle !_A, 1_A\rangle. \]
    For the other direction, we take $g : 1 \to B^A$, hence:
    \begin{align*}
        \ulcorner g^\sharp \urcorner &= ((\ev_{B,A} \circ \langle g\circ !_A,1_A\rangle) \circ \pi_{1,A}')^* &&(\textnormal{by the definitions}) \\
                                     &= (\ev_{B,A} \circ (\langle g\circ !_A, 1_A\rangle \circ \pi_{1,A}'))^* &&(\textnormal{by associativity law}) \\
                                     &= (\ev_{B,A} \circ \langle g\circ !_A \circ \pi_{1,A}', 1_A \circ \pi_{1,A}'\rangle)^* &&(\textnormal{by distributivity law}) \\
                                     &= (\ev_{B,A} \circ \langle g\circ \pi_{1,A}, 1_A \circ \pi_{1,A}'\rangle)^* &&(\textnormal{by $\spadesuit$}) \\
                                     &= (\ev_{B,A} \circ (g \times 1_A))^* &&(\textnormal{by the definition of $\times$}) \\
                                     &= g &&(\textnormal{by CCC2}),
    \end{align*}
    where, in the step marked by $\spadesuit$, we used the fact that $!_A \circ \pi_{1,A}' : 1 \times A \to 1$, since $1$ is the terminal object, then (CC1) gives $!_A \circ \pi_{1,A}' = !_{1 \times A} = \pi_{1,A}$.
\end{proof}

 
\section{Polynomial categories}
Now that we have defined the structure of a cartesian closed category, we want a way to construct a new category, from an existing one, adding a new arrow.
In order to do this, we first need to give a precise definition of a cartesian closed category generated by a graph.

We proceed in two steps, first we define the positive intuitionistic propositional calculus generated by a graph, then we define the cartesian closed category generated by a positive intuitionistic propositional calculus.

\begin{definition}[Freely generated positive intuitionistic calculus by a graph]
    Given a graph $\mathcal G$, the \textit{freely generated positive intuitionistic propositional calculus} generated by $\mathcal G$, written as $\mathcal P(\mathcal G)$,
    is the smallest positive intuitionistic propositional calculus whose formulas include the objects of $\mathcal G$ and whose proofs include the arrows of $\mathcal G$.
    More precisely, we can give a recursive definition of both the objects and the arrows of $\mathcal P(\mathcal G)$:
    \begin{itemize}
        \item the objects of $\mathcal P(\mathcal G)$ are the formulas generated by the following grammar:
        \[ A,B ::= \top \mid X \mid A \land B \mid A \Leftarrow B,\text{\footnotemark} \]
        \footnotetext{This is BNF-style notation, more precisely the variant known as Math-BNF (MBNF), commonly used to define mathematical syntax; its general form is $\bullet ::= \circ_1 \mid \cdots \mid \circ_n$ \cite{QuinlanWellsKamareddine2019}.
        Here and in all subsequent MBNF grammars for arrows, the metavariables (for example $f,g$) are sorted by source and target, and only well-typed instances of the displayed productions, according to the corresponding rules of inference, are admitted.}
        where $X$ is an object of $\mathcal G$;
        \item the arrows of $\mathcal P(\mathcal G)$ are the proofs generated by the following grammar:
        \[ f,g ::= h \mid 1_A \mid !_A \mid f \circ g \mid \pi_{A,B} \mid \pi_{A,B}' \mid \langle f,g\rangle \mid f^* \mid \ev_{A,B},\]
        where $f$ and $g$ are arrows of $\mathcal P(\mathcal G)$, $A$ and $B$ are objects of $\mathcal P(\mathcal G)$, and $h$ is an arrow of $\mathcal G$.
    \end{itemize} 
\end{definition}

Now that we have the positive intuitionistic calculus generated by $\mathcal G$, we can define the cartesian closed category generated by it. We'll do that in two steps, beginning with the following lemma.

\begin{lemma}\label{lem:ccc-quotient-existence}
Given a graph $\mathcal G$, there exists an equivalence relation $\sim$ on the
arrows of $\mathcal P(\mathcal G)$, compatible with composition,
pairing, currying, and containing the category, cartesian and
cartesian closed axioms. Consequently,
$\mathcal P(\mathcal G)/{\sim}$ is a cartesian closed category.
\end{lemma}

\begin{proof}
    We already said in \Cref{rem:category-generated-by-deductive-system} that if we take the equivalence relation of having the same source and target, 
    this is compatible with composition, and it contains the category axioms; if we check that: it is compatible with the operations (inference rules) of
    pairing, currying, and it contains the cartesian and cartesian closed axioms, then we can conclude that the quotient is a cartesian closed category.

    Let's check point by point:
    \begin{itemize}
        \item for the pairing, if we take $f,f' : A \to B$, $g,g' : A \to C$, then:
        \[ \langle f,g\rangle : A \to B \land C, \qquad \langle f',g'\rangle : A \to B \land C, \]
        therefore $\langle f,g\rangle$ and $\langle f',g'\rangle$ have the same source and target, so they are equivalent;
        \item for the currying, if we take $h,h' : A \land B \to C$, then:
        \[ h^* : A \to C \Leftarrow B, \qquad h'^* : A \to C \Leftarrow B, \]
        therefore $h^* \sim (h')^*$;
        \item for the cartesian axioms, we have:
        \begin{itemize}[label=$\ast$]
            \item for (CC1), given $f : A \to \top$, we have $f \sim !_A$;
            \item for (CC2a), given $f : C \to A$, $g : C \to B$, we have $\pi_{A,B} \circ \langle f,g\rangle \sim f$, by their definition;
            \item similarly for (CC2b);
            \item for (CC2c), given $h : C \to A \land B$, we have $\langle \pi_{A,B} \circ h, \pi_{A,B}' \circ h\rangle : C \to A \land B$, therefore it's equivalent to $h$;
        \end{itemize}
        \item for the cartesian closed axioms, we have:
        \begin{itemize}[label=$\ast$]
            \item for (CCC1), given $h : A \land B \to C$, we have $\ev_{C,B} \circ (h^* \wedge 1_B) \sim h$;
            \item for (CCC2), given $k : A \to C \Leftarrow B$, we have $(\ev_{C,B} \circ (k \wedge 1_B))^* \sim k$.
        \end{itemize}
    \end{itemize}
    Since the equivalence relation of having the same source and target is compatible with composition, pairing, currying, and contains the category, cartesian and cartesian closed axioms, we can conclude that the quotient $\mathcal P(\mathcal G)/{\sim}$ is a cartesian closed category.
\end{proof}

We can now give the definition of a cartesian closed category generated by a graph, which will be well-defined thanks to the previous lemma.

\begin{definition}[Cartesian closed category generated by a graph]
    Given a graph $\mathcal G$, the \textit{cartesian closed category freely generated} by $\mathcal G$, written as $\mathcal F(\mathcal G)$,
    is the quotient of $\mathcal P(\mathcal G)$ by the smallest equivalence relation on parallel arrows that contains the category, cartesian and cartesian closed axioms and is compatible with composition, pairing and currying.
\end{definition}

The cartesian closed category freely generated by a graph has an important universal property, which will be important later.
Before stating and proving it, let's define the notion of functor between cartesian closed categories.

\begin{definition}[Cartesian closed functor]
    Given two cartesian closed categories $\mathcal C$ and $\mathcal D$, a \textit{cartesian closed functor} $F : \mathcal C \to \mathcal D$ is a functor with the following properties:
    \begin{enumerate}[label=(\textit{\roman*})]
        \item $F(1_\mathcal C) = 1_\mathcal D$;
        \item $F(A \times B) = F(A) \times F(B)$;
        \item $F(A^B) = F(A)^{F(B)}$;
        \item $F(!_A) = !_{F(A)}$;
        \item $F(\pi_{A,B}) = \pi_{F(A),F(B)}$;
        \item $F(\pi_{A,B}') = \pi_{F(A),F(B)}'$;
        \item $F(\langle f,g\rangle) = \langle F(f), F(g)\rangle$;
        \item $F(f^*) = F(f)^*$;
        \item $F(\ev_{A,B}) = \ev_{F(A),F(B)}$.
    \end{enumerate}
\end{definition}

\begin{notation}
    We denote by $\mathsf{CCC}$ the category of cartesian closed categories (whose morphisms are cartesian closed functors) and by $\mathsf{Grph}$ the category of graphs (whose morphisms are maps that commute with the functions $s$ and $t$).
    Moreover, we will denote by $\mathcal U : \mathsf{CCC} \to \mathsf{Grph}$ the forgetful functor that sends a cartesian closed category to its underlying graph.
\end{notation}

\begin{definition}[Canonical inclusion into the free CCC underlying graph]
    Given a graph $\mathcal G$, we define the following morphism of graphs:
    \[ H_\mathcal G : \mathcal G \to \mathcal U(\mathcal F(\mathcal G)) : X \mapsto X, f \mapsto [f],\]
    that is each object, $X$ of $\mathcal G$, is sent to itself in $\mathcal F(\mathcal G)$ and
    each arrow $f : X \to Y$ of $\mathcal G$ is sent to its equivalence class $[f]$ in $\mathcal U(\mathcal F(\mathcal G))$. That is, the same 
    equivalence class of $f$ as in $\mathcal F(\mathcal G)$. We're simply viewing it in the underlying graph structure and forgetting the cartesian closed structure.
\end{definition}

We can now state the universal property of the cartesian closed category freely generated by a graph.

\begin{theorem}[Universal property of CCCs generated by a graph]
    Given a cartesian closed category $\mathcal C$, a graph $\mathcal G$ and a morphism of graphs $F : \mathcal G \to \mathcal U(\mathcal C)$,
    there is a unique cartesian closed functor $\overline F : \mathcal F(\mathcal G) \to \mathcal C$ such that the following diagram commutes:
    % https://q.uiver.app/#q=WzAsMyxbMCwxLCJcXG1hdGhjYWwgVShcXG1hdGhjYWwgRihcXG1hdGhjYWwgRykpIl0sWzAsMCwiXFxtYXRoY2FsIEciXSxbMSwwLCJcXG1hdGhjYWwgVShcXG1hdGhjYWwgQykiXSxbMCwyLCJcXG1hdGhjYWwgVShcXG92ZXJsaW5lIEYpIiwyLHsic3R5bGUiOnsiYm9keSI6eyJuYW1lIjoiZGFzaGVkIn19fV0sWzEsMiwiRiJdLFsxLDAsIkhfe1xcbWF0aGNhbCBHfSIsMl1d
    \[\begin{tikzcd}
        {\mathcal G} & {\mathcal U(\mathcal C)} \\
        {\mathcal U(\mathcal F(\mathcal G))}
        \arrow["F", from=1-1, to=1-2]
        \arrow["{H_{\mathcal G}}"', from=1-1, to=2-1]
        \arrow["{\mathcal U(\overline F)}"', dashed, from=2-1, to=1-2]
    \end{tikzcd}.\]
\end{theorem}

\begin{proof}
    We need to uniquely define a cartesian closed functor $\overline F : \mathcal F(\mathcal G) \to \mathcal C$, we know that the objects of $\mathcal F(\mathcal G)$ coincide with the objects of the positive intuitionistic calculus $\mathcal P(\mathcal G)$,
    which, using categorical notation, are generated by the objects of $\mathcal G$ through the grammar:
    \[ A,B ::= 1 \mid X \mid A \times B \mid A^B, \]
    where $X \in \ob(\mathcal G)$. Since we want $\mathcal U(\overline F) \circ H_\mathcal G = F$, for every $X \in \ob(\mathcal G)$, we must have
    \[ F(X) = \mathcal U(\overline F)\circ (H_\mathcal G)(X) = \mathcal U(\overline F)(X)  \qquad (\text{by definition of $H_{\mathcal G}$})\]
    and because $\mathcal U(\overline F)$ coincides with $\overline F$ on objects, by definition, we have that $\overline F(X) = F(X)$, $\forall X \in \ob(\mathcal G)$.
    This definition of $\overline F$ can be recursively extended to the objects of $\mathcal P(\mathcal G)$ by imposing the cartesian closed structure, therefore $\overline F$ is uniquely defined on the objects of $\mathcal F(\mathcal G)$\footnote{To be completely rigorous, we should invoke the principles of recursion on the grammar, to define $\overline F$, and induction on the grammar, to prove its uniqueness.}.\par
    By the same reasoning, we can recursively define $\overline F$ on the arrows of $\mathcal P(\mathcal G)$ - the positive intuitionistic calculus generated by $\mathcal G$, so it's cartesian closed, since they are generated by the grammar:
    \[ f,g ::= h \mid 1_A \mid !_A \mid f \circ g \mid \pi_{A,B} \mid \pi_{A,B}' \mid \langle f,g\rangle \mid f^* \mid \ev_{A,B},\]
    where $h \in \mor(\mathcal G)$, $A$ and $B$ are objects of $\mathcal F(\mathcal G)$, and $f, g \in \mor(\mathcal F(\mathcal G))$.\par
    We can then extend, by checking its compatibility with the equivalence relation, the definition of $\mathcal F(\mathcal G)$ to the arrows of $\mathcal P(\mathcal G)$, doing this we get a uniquely well-defined map on the arrows too.\\
    Since the equivalence relation used to define $\mathcal F(\mathcal G)$ is the smallest equivalence relation on parallel arrows compatible with composition, pairing and currying, and containing the category, cartesian and cartesian closed axioms, 
    we can prove that the relation of having the same image under $\overline F$ is an equivalence relation compatible with all the mentioned operations and containing the axioms, hence it contains the equivalence relation used to define $\mathcal F(\mathcal G)$,
    so we would have that $[f] = [g]$ implies $\overline F(f) = \overline F(g)$ (because it's true for a relation that contains ours).\\
    We just need to prove that $\sim_{\overline F}$, where $f \sim_{\overline F} g$ if and only if $\overline F(f) = \overline F(g)$ and $f,g \in \Hom(A,B)$ for some $A$ and $B$, has all the properties listed above:
    \begin{itemize}
        \item it is obviously an equivalence relation, since equality is an equivalence relation;
        \item it is compatible with composition, since if $\overline F(f) = \overline F(g)$ and $\overline F(f') = \overline F(g')$ then, by the functoriality of $\overline F$:
        \[ \overline F(f \circ f') = \overline F(f) \circ \overline F(f') = \overline F(g) \circ \overline F(g') = \overline F(g \circ g');\]
        \item it is compatible with pairing, in the sense that if $\overline F(f) = \overline F(f')$ and $\overline F(g) = \overline F(g')$ then, by the cartesian closedness of $\overline F$:
        \[ \overline F(\langle f,g\rangle) = \langle \overline F(f), \overline F(g)\rangle = \langle \overline F(f'), \overline F(g')\rangle = \overline F(\langle f',g'\rangle);\]
        \item it is compatible with currying, with the same argument, i.e. if $\overline F(h) = \overline F(h')$ then:
        \[ \overline F(h^*) = \overline F(h)^* = \overline F(h')^* = \overline F((h')^*);\]
        \item it contains the category axioms:
        \begin{itemize}[label=$\ast$]
            \item $1_Bf \sim_{\overline F} f$, since $\overline F(1_Bf) = \overline F(1_B)\overline F(f) = 1_{\overline F(B)}\overline F(f) = \overline F(f)$, and similarly for the other identity law;
            \item $(h \circ g) \circ f \sim_{\overline F} h \circ (g \circ f)$, also by the functoriality of $\overline F$;
        \end{itemize}
        \item it contains the cartesian axioms:
        \begin{itemize}[label=$\ast$]
            \item for (CC1), given $f : A \to 1$, we have $\overline F(f) : \overline F(A) \to 1$, therefore $\overline F(f) = !_{\overline F(A)}$ in $\mathcal C$, and $\overline F(!_A) = !_{\overline F(A)}$, by the cartesian closedness;
            \item for (CC2a), given $f : C \to A$, $g : C \to B$, we have:
            \begin{align*}
                \overline F(\pi_{A,B} \circ \langle f,g\rangle) &= \overline F(\pi_{A,B}) \circ \overline F(\langle f,g\rangle) &&(\textnormal{by functoriality}) \\
                &= \pi_{\overline F(A),\overline F(B)} \circ \langle \overline F(f), \overline F(g)\rangle &&(\textnormal{by cartesian closedness}) \\
                &= \overline F(f) &&(\textnormal{by CC2a in $\mathcal C$});
            \end{align*}
            \item similarly for (CC2b);
            \item for (CC2c), given $h : C \to A \times B$, we have:
            \begin{align*}
                \overline F(\langle \pi_{A,B} \circ h, \pi_{A,B}' \circ h\rangle) &= \langle \overline F(\pi_{A,B}) \circ \overline F(h), \overline F(\pi_{A,B}') \circ \overline F(h)\rangle \\
                                                                                  &= \langle \pi_{\overline F(A),\overline F(B)} \circ \overline F(h), \pi_{\overline F(A),\overline F(B)}' \circ \overline F(h)\rangle \\
                                                                                  &= \overline F(h) &&(\textnormal{by CC2c in $\mathcal C$});
            \end{align*}
        \end{itemize}
        \item it contains the cartesian closed axioms:
        \begin{itemize}[label=$\ast$]
            \item for (CCC1), given $h : A \times B \to C$, we have:
            \begin{align*}
                \overline F(\ev_{C,B} \circ (h^* \times 1_B)) &= \overline F(\ev_{C,B}) \circ \overline F(h^* \times 1_B) \\
                                                              &= \ev_{\overline F(C),\overline F(B)} \circ (\overline F(h)^* \times 1_{\overline F(B)}) \\
                                                              &= \overline F(h) &&(\textnormal{by CCC1 in $\mathcal C$}).
            \end{align*}
            \item for (CCC2), given $k : A \to C^B$, we have:
            \begin{align*}
                \overline F((\ev_{C,B} \circ (k \times 1_B))^*) &= (\overline F(\ev_{C,B} \circ (k \times 1_B)))^* \\
                                                                 &= ((\ev_{\overline F(C),\overline F(B)}) \circ (\overline F(k) \times 1_{\overline F(B)}))^* \\
                                                                 &= \overline F(k) &&(\textnormal{by CCC2 in $\mathcal C$}).
            \end{align*}
        \end{itemize}
    \end{itemize}
\end{proof}

\begin{remark}[$\mathcal F$ is a functor]
    From the universal property just seen we can easily deduce that $\mathcal F : \mathsf{Grph} \to \mathsf{CCC}$ is a functor. We already know how it associates a cartesian closed category $\mathcal F(\mathcal G)$ to each graph $\mathcal G$, 
    now we can define $\mathcal F$ on graph morphisms. Given $\varphi : \mathcal G \to \mathcal H$, we can define $\widetilde{\varphi} : \mathcal G \to \mathcal U(\mathcal F(\mathcal H))$ as $\widetilde{\varphi} := H_\mathcal H \circ \varphi$, then, by the universal property, we have the following commutative diagram:
    % https://q.uiver.app/#q=WzAsNCxbMCwxLCJcXG1hdGhjYWwgRyJdLFswLDAsIlxcbWF0aGNhbCBIIl0sWzEsMSwiXFxtYXRoY2FsIFUoXFxtYXRoY2FsIEYoXFxtYXRoY2FsIEgpKSJdLFswLDIsIlxcbWF0aGNhbCBVKFxcbWF0aGNhbCBGKFxcbWF0aGNhbCBHKSkiXSxbMCwxLCJcXHZhcnBoaSJdLFsxLDIsIkhfe1xcbWF0aGNhbCBIfSJdLFswLDIsIlxcd2lkZXRpbGRle1xcdmFycGhpfSIsMl0sWzAsMywiSF97XFxtYXRoY2FsIEd9IiwyXSxbMywyLCJcXG1hdGhjYWx7VX0oXFxvdmVybGluZXtcXHdpZGV0aWxkZXtcXHZhcnBoaX19KSIsMix7InN0eWxlIjp7ImJvZHkiOnsibmFtZSI6ImRhc2hlZCJ9fX1dXQ==
    \[\begin{tikzcd}
        {\mathcal H} & \\
        {\mathcal G} & {\mathcal U(\mathcal F(\mathcal H))} \\
        {\mathcal U(\mathcal F(\mathcal G))}
        \arrow["{H_{\mathcal H}}", from=1-1, to=2-2]
        \arrow["\varphi", from=2-1, to=1-1]
        \arrow["{\widetilde{\varphi}}"', from=2-1, to=2-2]
        \arrow["{H_{\mathcal G}}"', from=2-1, to=3-1]
        \arrow["{\mathcal{U}(\overline{\widetilde{\varphi}})}"', dashed, from=3-1, to=2-2]
    \end{tikzcd},\]
where $\overline{\widetilde{\varphi}} : \mathcal F(\mathcal G) \to \mathcal F(\mathcal H)$ is the unique cartesian closed functor given by the universal property above, so we can define $\mathcal F(\varphi) := \overline{\widetilde{\varphi}} = \overline{H_\mathcal H \circ \varphi}$. One can 
easily check that $\mathcal F$ preserves identities and composition using the functoriality of $\mathcal U$ and the uniqueness given by the universal property. 
\end{remark}

\begin{remark}
    This universal property can be equivalently expressed as the following adjunction:
    \[ \Hom_{\mathsf{CCC}}(\mathcal F(\mathcal G), \mathcal C) \overset{1:1}{\longleftrightarrow} \Hom_{\mathsf{Grph}}(\mathcal G, \mathcal U(\mathcal C)), \]
    natural in both $\mathcal G$ and $\mathcal C$, i.e., compatible with precomposition in $\mathcal G$\footnote{Which means precomposing with a graph morphism into $\mathcal G$ and then extending it with the universal property is the same as extending and then composing. The same applies to the postcomposition with cartesian closed functors out of $\mathcal C$.} and postcomposition in $\mathcal C$.
    Therefore the free-CCC functor, $\mathcal F : \mathsf{Grph} \to \mathsf{CCC}$, is left adjoint to the forgetful functor $\mathcal U : \mathsf{CCC} \to \mathsf{Grph}$; or equivalently $\mathcal U$ is right adjoint to $\mathcal F$ (so we can write $\mathcal F \dashv \mathcal U$).
\end{remark}

Finally, we now define the core notion of this section, that is the notion of polynomial category over a cartesian closed category. The question is:
given two objects $A_0$ and $A$ of a (cartesian, cartesian closed) category $\mathcal C$, how does one adjoin an \textit{indeterminate} arrow $x : A_0 \to A$ to $\mathcal C$?

\begin{definition}[Polynomial category - v.1]
    Given a (cartesian, cartesian closed) category $\mathcal C$, two objects $A_0$ and $A$ of $\mathcal C$, and an indeterminate arrow $x : A_0 \to A$, we define the \textit{polynomial category} $\mathcal C[x]$ 
    as the cartesian closed category freely generated by the graph $\mathcal U(\mathcal C)$ with the additional arrow $x : A_0 \to A$, quotiented in order to preserve the relations between the arrows of $\mathcal C$ and make the natural immersion a functor.
\end{definition}

\begin{notation}[Polynomial arrows]
    We will denote by $\varphi(x), \psi(x), \ldots$ the arrows of $\mathcal C[x]$ that are generated by the indeterminate arrow $x : A_0 \to A$. These arrows 
    can be thought of as proofs under the \textit{assumption} $x : A_0 \to A$.
\end{notation}

The previous definition is not precise, but it conveys the general idea. It can be made precise by following a similar approach to the one used to define the cartesian closed category freely generated by a graph.
This provides us with an explicit description of the arrows of $\mathcal C[x]$ and of the equivalence relation used to define it, which will be useful later.

\begin{definition}[Polynomial category - v.2]
    Given a (cartesian, cartesian closed) category $\mathcal C$, two objects $A_0$ and $A$ of $\mathcal C$, and an indeterminate arrow $x : A_0 \to A$, we define the \textit{polynomial category} $\mathcal C[x]$
    in two steps:
    \begin{itemize}
        \item first we define the positive intuitionistic propositional calculus, this time, since $\mathcal C$ is already a cartesian closed category, we can directly take its objects $\ob(\mathcal C)$ as the objects of the calculus. Moreover, the tautological object, conjunction and implication functions are those of $\mathcal C$. For the arrows,
        we take the ones generated by the following grammar:
        \[ f,g ::= h \mid x \mid f \circ g \mid \langle f,g\rangle \mid f^*, \]
        where $A, B \in \ob(\mathcal C)$, $h \in \mor(\mathcal C)$ and $x : A_0 \to A$ is the indeterminate arrow. In particular, the identity, terminal, projection and evaluation arrows of $\mathcal C$ are included among the arrows $h$;
        \item now we need to impose the equivalence relation $\sim$ on the arrows so that the quotient $\mathcal C[x]$ is a (cartesian, cartesian closed) category and the immersion
        \[ \iota_x : \mathcal C \rightarrow \mathcal C[x] : X \mapsto X, f \mapsto [f], \]
        is a (cartesian, cartesian closed) functor. We choose the smallest equivalence relation on parallel arrows satisfying the following properties, and we denote it by $\underset{x}{=}$:
        \begin{itemize}[label=$\ast$]
            \item if $\varphi(x) \underset{x}{=} \varphi'(x)$ and $\psi(x) \underset{x}{=} \psi'(x)$, then $\varphi(x) \circ \psi(x) \underset{x}{=} \varphi'(x) \circ \psi'(x)$;
            \item if $\varphi(x) \underset{x}{=} \varphi'(x)$ and $\psi(x) \underset{x}{=} \psi'(x)$, then $\langle \varphi(x), \psi(x)\rangle \underset{x}{=} \langle \varphi'(x), \psi'(x)\rangle$;
            \item if $\varphi(x) \underset{x}{=} \varphi'(x)$, then $\varphi(x)^* \underset{x}{=} \varphi'(x)^*$;
            \item it contains the category, cartesian and cartesian closed axioms;
            \item $g \circ f = h$ in $\mathcal C$ implies $g \circ f \underset{x}{=} h$ in $\mathcal C[x]$ (the composition in $\mathcal C[x]$ extends the one in $\mathcal C$);
            \item for each arrow $f : C \to A$ and $g : C \to B$ of $\mathcal C$, if $\langle f,g\rangle = h$, then $\langle f,g\rangle \underset{x}{=} h$ in $\mathcal C[x]$ (the pairing in $\mathcal C[x]$ extends the one in $\mathcal C$);
            \item if $f^* = h$, then $f^* \underset{x}{=} h$ (the currying in $\mathcal C[x]$ extends the one in $\mathcal C$);
        \end{itemize}
    \end{itemize}
\end{definition}

\begin{remark}[Well-posedness of the definition v.2]
    By \Cref{lem:ccc-quotient-existence}, we already know that the relation of having the same source and target is compatible with composition, pairing, currying, and contains the category, cartesian and cartesian closed axioms.
    One can easily check that this equivalence relation also contains the remaining properties above, therefore the quotient is a (cartesian, cartesian closed) category and the immersion $\iota_x : \mathcal C \rightarrow \mathcal C[x]$ is a (cartesian, cartesian closed) functor,
    so the definition above is well-posed.
\end{remark}

\begin{notation}[Canonical functor $\mathcal C \to \mathcal C{[x]}$]
    From now on we will denote by $\iota_x : \mathcal C \rightarrow \mathcal C[x]$ the immersion of $\mathcal C$ into $\mathcal C[x]$, which is a (cartesian, cartesian closed) functor.
    When it is clear from the context, we will sometimes omit the subscript $x$ and just write $\iota : \mathcal C \rightarrow \mathcal C[x]$.
\end{notation}

The polynomial category $\mathcal C[x]$ has an important universal property, which is the following.

\begin{theorem}[Universal property of polynomial categories]
    Given a (cartesian, cartesian closed) category $\mathcal C$, two objects $A_0$ and $A$ of $\mathcal C$, an indeterminate arrow $x : A_0 \to A$, a (cartesian, cartesian closed) category $\mathcal D$, a (cartesian, cartesian closed) functor $F : \mathcal C \to \mathcal D$ and an arrow $a : F(A_0) \to F(A)$ in $\mathcal D$,
    there is a unique (cartesian, cartesian closed) functor $\widehat F : \mathcal C[x] \to \mathcal D$ such that $\widehat F([x]) = a$ and the following diagram commutes:
    % https://q.uiver.app/#q=WzAsMyxbMCwxLCJcXG1hdGhjYWwgQ1t4XSJdLFswLDAsIlxcbWF0aGNhbCBDIl0sWzEsMCwiXFxtYXRoY2FsIEQiXSxbMSwwLCJcXGlvdGFfeCIsMl0sWzEsMiwiRiJdLFswLDIsIlxcd2lkZWhhdCBGIiwyLHsic3R5bGUiOnsiYm9keSI6eyJuYW1lIjoiZGFzaGVkIn19fV1d
    \[\begin{tikzcd}
        {\mathcal C} & {\mathcal D} \\
        {\mathcal C[x]}
        \arrow["F", from=1-1, to=1-2]
        \arrow["{\iota_x}"', from=1-1, to=2-1]
        \arrow["{\widehat F}"', dashed, from=2-1, to=1-2]
    \end{tikzcd}.\]
\end{theorem}

\begin{proof}
    Let's prove the theorem in the case that everything is cartesian closed, the other cases follow. We need to uniquely define the cartesian closed functor $\widehat F : \mathcal C[x] \to \mathcal D$,
    since the objects of $\mathcal C[x]$ coincide with the objects of $\mathcal C$, we must have $\widehat F(A) := F(A)$ for every object $A$ of $\mathcal C$, therefore $\widehat F$ is cartesian closed on objects, makes the diagram commute, and is uniquely determined on objects.\\
    The arrows of $\mathcal C[x]$ are the equivalence classes of the arrows generated by the grammar:
    \[ f,g ::= h \mid x \mid f \circ g \mid \langle f,g\rangle \mid f^*,\]
    where $h \in \mor(\mathcal C)$ and $A,B \in \ob(\mathcal C)$. So a proof $\varphi(x)$ in the positive calculus has necessarily one of the following forms:
    \[ f,x,\psi(x) \circ \varphi(x), \langle \psi(x), \varphi(x)\rangle, \varphi(x)^*, \]
    with $f \in \mor(\mathcal C)$, hence we define:
    \begin{align*}
        \widehat F([f]) &:= F(f), &&(\textnormal{hence the diagram commutes})\\
        \widehat F([x]) &:= a, \\
        \widehat F([\psi(x) \circ \varphi(x)]) &:= \widehat F([\psi(x)]) \circ \widehat F([\varphi(x)]), \\
        \widehat F([\langle \psi(x), \varphi(x)\rangle]) &:= \langle \widehat F([\psi(x)]), \widehat F([\varphi(x)])\rangle, \\
        \widehat F([\varphi(x)^*]) &:= \widehat F([\varphi(x)])^*.
    \end{align*}
    If we check that $\widehat F$ is well-defined on the equivalence classes of arrows, then it will automatically be a cartesian closed functor, it will make the diagram commute, and it will be uniquely determined by the definition above\footnote{As we remarked before, to be completely rigorous, we should invoke recursion and induction on the grammar.}.\par
    The idea for proving the well-definedness is the same as in the proof of the other universal property already seen, we define the equivalence relation on the arrows, $f \sim_{\widehat F} g$ if and only if $\widehat F(f) = \widehat F(g)$ and $f,g \in \Hom(A,B)$ for some $A$ and $B$, and we prove that such relation is compatible with composition, pairing and currying, contains the category, cartesian and cartesian closed axioms, and satisfies the extension properties listed in the definition,
    so it must contain $\underset{x}{=}$, and this will imply that if $[f] = [g]$ then $\widehat F(f) = \widehat F(g)$. All the checks are similar to the ones in the other proof and follow the scheme in the definition v.2, so we omit them here. 
\end{proof}

\section{Functional completeness and Lambek's normal form}
The aim of this section is to provide a normal form for the arrows of the polynomial category $\mathcal C[x]$ of a cartesian closed category $\mathcal C$.
Such normal form should be unique and is based on the same idea as the normal form of polynomials in $R[x]$, that is, for a commutative unital ring $R$,
every polynomial can be uniquely expressed in the form $a_n x^n + a_{n-1} x^{n-1} + \ldots + a_1 x + a_0$, with $a_i \in R$, $n \in \NN$.\par
We'll try to do something similar for the arrows of $\mathcal C[x]$ and obtain \textit{Lambek's normal form}. In order to do this we first need to introduce a couple 
of results, each a stronger version of the previous one. These results go under the name of \textit{deduction theorems} or, more broadly, \textit{functional completeness} theorems.

\begin{theorem}[Deduction theorem for positive intuitionistic calculi]
    Given a positive intuitionistic propositional calculus $\mathcal P$, an object $A$ of $\mathcal P$ and an indeterminate arrow $x : \top \to A$,
    for every arrow $\varphi(x) : B \to C$ of $\mathcal P\langle x\rangle$, there exists an arrow:
    \[ \kappa_{x \in A}\varphi(x) : A \land B \to C, \]
    that belongs to $\mor(\mathcal P)$.
\end{theorem}

By $\mathcal P\langle x\rangle$ we mean the formal polynomial positive intuitionistic calculus over $\mathcal P$, which is defined as follows.

\begin{definition}
    Given a positive intuitionistic propositional calculus $\mathcal P$, an object $A$ of $\mathcal P$ and an indeterminate arrow $x : \top \to A$, we denote by
    $\mathcal P\langle x\rangle$ the \textit{formal polynomial positive intuitionistic calculus over $\mathcal P$}. Its objects, tautological object, conjunction and implication are those of
    $\mathcal P$, while its arrows are generated by the grammar
    \[ \varphi,\psi ::= h \mid x \mid \psi \circ \varphi \mid \langle \varphi,\psi\rangle \mid \varphi^*, \]
    where $h \in \mor(\mathcal P)$ and $A \in \ob(\mathcal P)$. In particular, the identity, terminal, projection and evaluation arrows $1_B$, $!_B$, $\pi_{B,C}$, $\pi_{B,C}'$ and $\textnormal{ev}_{B,C}$ of $\mathcal P$ are the same of $\mathcal P\langle x \rangle$, since they are included in the $h$ above.
    We do not impose any equivalence relation, so, for example, if $g \circ f = h$ in $\mathcal P$, then the term $h$ and the formal composition $g \circ f$ are distinct arrows of $\mathcal P\langle x\rangle$.
\end{definition}

\begin{example}
    We know that a cartesian closed category $\mathcal C$ is a positive intuitionistic propositional calculus, so if we take an object $A$ and $x : 1 \to A$, then,
    by the definitions, we have $\mathcal C[x] = \mathcal C\langle x\rangle / \underset{x}{=}$.
\end{example}

\begin{remark}[Other versions of the deduction theorem]
    The theorem above can also be stated and proved for conjunction, positive intuitionistic propositional, intuitionistic propositional and classical calculi, using the corresponding formal polynomial term calculi.
    The idea of the proof is the same as the one we are about to show, with the appropriate cases added to or removed from the induction.\newline
    We prove the theorem for a positive intuitionistic propositional calculus because it is the syntactic construction needed for the next result on cartesian closed categories.
\end{remark}

We now give a constructive proof of the deduction theorem stated above.

\begin{proof}
    By the definition of $\mathcal P\langle x\rangle$, every formal arrow $\varphi(x) : B \to C$ has one of the following forms:
    \[ g,\quad x,\quad \psi(x) \circ \chi(x),\quad \langle \psi(x),\chi(x)\rangle,\quad \psi(x)^*, \]
    where $g \in \mor(\mathcal P)$ and $x : \top \to A$ is the indeterminate arrow. We proceed by induction on the \textit{length} $\ell$ of the term, that is defined as follows:
    \begin{itemize}
        \item if $\varphi(x) = g$ or $\varphi(x) = x$, then $\ell(\varphi(x)) := 1$;
        \item if $\varphi(x) = \psi(x) \circ \chi(x)$, then $\ell(\varphi(x)) := \ell(\psi(x)) + \ell(\chi(x)) + 1$;
        \item if $\varphi(x) = \langle \psi(x), \chi(x)\rangle$, then $\ell(\varphi(x)) := \ell(\psi(x)) + \ell(\chi(x)) + 1$;
        \item if $\varphi(x) = \psi(x)^*$, then $\ell(\varphi(x)) := \ell(\psi(x)) + 1$.
    \end{itemize} 
    For each of the five cases, we define the associated arrow recursively, using the induction hypothesis on shorter terms:
    \begin{itemize}
        \item if $\varphi(x) = g$, where $g : B \to C$, then we define $\kappa_{x \in A}g := g \circ \pi_{A,B}'$;
        \item if $\varphi(x) = x$, then $B = \top$ and $C = A$, so we define $\kappa_{x \in A}x := \pi_{A,\top}$;
        \item if $\varphi(x) = \psi(x) \circ \chi(x)$, where $\chi(x) : B \to D$ and $\psi(x) : D \to C$, then $\kappa_{x \in A}\chi(x) : A \land B \to D$ and $\kappa_{x \in A}\psi(x) : A \land D \to C$, so we define
        \[ \kappa_{x \in A} (\psi(x) \circ \chi(x)) := \kappa_{x \in A} \psi(x) \circ \langle \pi_{A,B}, \kappa_{x \in A} \chi(x)\rangle : A \land B \to C; \]
        \item if $\varphi(x) = \langle \psi(x), \chi(x)\rangle$, where $\psi(x) : B \to D$, $\chi(x) : B \to E$ and $C = D \land E$, then we define
        \[ \kappa_{x \in A} \langle \psi(x), \chi(x)\rangle := \langle \kappa_{x \in A} \psi(x), \kappa_{x \in A} \chi(x)\rangle : A \land B \to C; \]
        \item if $\varphi(x) = \psi(x)^*$, where $\psi(x) : B \land D \to E$ and $C = E \Leftarrow D$, then $\kappa_{x \in A}\psi(x) : A \land (B \land D) \to E$. We define
        \[ \kappa_{x \in A}\big(\psi(x)^*\big) := \big(\kappa_{x \in A}\psi(x) \circ \alpha_{A,B,D}\big)^* : A \land B \to C, \]
        where $\alpha_{A,B,D} : (A \land B) \land D \to A \land (B \land D)$ denotes the \hyperref[ex:associativity-of-conjunction]{associativity arrow}.
    \end{itemize}
    In every case the displayed expression is an arrow of $\mathcal P$ of the required type. The recursive construction therefore yields
    $\kappa_{x \in A}\varphi(x) : A \land B \to C$ for every formal term $\varphi(x) : B \to C$.
\end{proof}

\begin{notation}
    From now on we will denote by $\kappa_{x \in A}\varphi(x)$ the arrow in $\mathcal P$ associated to $\varphi(x)$ by the deduction theorem above.
    The subscript $x \in A$ indicates that $x$ is of \textit{type} $A$, i.e.,\ $x : \top \to A$.
\end{notation}

\begin{remark}[Classical deduction theorem]
    [PLACEHOLDER]
\end{remark}

\begin{theorem}[Functional completeness]
    Given a cartesian closed category $\mathcal C$, an object $A$ of $\mathcal C$ and an indeterminate arrow $x : 1 \to A$, for every arrow $\varphi(x) : B \to C$ of $\mathcal C[x]$
    there is a unique arrow $f : A \times B \to C$ in $\mathcal C$ such that:
    \[ f \circ \langle x \circ !_B, 1_B\rangle \underset{x}{=} \varphi(x). \]
\end{theorem}

\begin{proof}
    The idea is that $f = \kappa_{x \in A}\varphi(x)$ works, but we first need to check that $\kappa_{x \in A}(\cdot)$ is well-defined on the equivalence classes of arrows of $\mathcal C[x]$. The uniqueness will follow easily after that.\newline
    \begin{itemize}
        \item[$\boxed{\begin{array}{c}\textnormal{well-}\\\textnormal{definedness}\end{array}}$] As usual, we define the equivalence relation on parallel arrows of $\mathcal C\langle x\rangle$
        \[ \varphi(x) \sim_{\kappa} \psi(x) \overset{\textnormal{def}}{\iff} \kappa_{x \in A}\varphi(x) = \kappa_{x \in A}\psi(x), \]
        and we check that it is compatible with composition, pairing and currying, contains the category, cartesian and cartesian closed axioms, and satisfies the extension properties listed in the definition of $\underset{x}{=}$. We omit the details.
        \item[$\boxed{\begin{array}{c}\kappa_{x \in A}\varphi(x)\\\textnormal{works}\end{array}}$] Let's check by induction on the length of $\varphi(x)$ that $f = \kappa_{x \in A}\varphi(x)$ satisfies the equality:
        \begin{itemize}
            \item if $\varphi(x) = g$, where $g : B \to C$, then $\kappa_{x \in A}g = g \circ \pi_{A,B}'$ and:
            \begin{align*}
                \kappa_{x \in A}g \circ \langle x \circ !_B, 1_B\rangle &\underset{x}{=} (g \circ \pi_{A,B}') \circ \langle x \circ !_B, 1_B\rangle \\
                                                                        &\underset{x}{=} g \circ 1_B \underset{x}{=} g;
            \end{align*}
            \item if $\varphi(x) = x$, then $\kappa_{x \in A}x = \pi_{A,1}$ and:
            \begin{align*}
                \kappa_{x \in A}x \circ \langle x \circ !_1, 1_1\rangle &\underset{x}{=} \pi_{A,1} \circ \langle x \circ !_1, 1_1\rangle \\
                                                                        &\underset{x}{=} x \circ !_1 \underset{x}{=} x;
            \end{align*}
            \item if $\varphi(x) = \psi(x) \circ \chi(x)$, with $\chi(x) : B \to D$ and $\psi(x) : D \to C$, then:
            \begin{align*}
                \kappa_{x \in A}(\psi(x) \circ \chi(x)) \circ \langle x \circ !_B, 1_B\rangle &\underset{x}{=} \kappa_{x \in A}(\psi(x)) \circ \langle \pi_{A,B}, \kappa_{x \in A}(\chi(x))\rangle \circ \langle x \circ !_B, 1_B\rangle \\
                                                                                              &\underset{x}{=} \kappa_{x \in A}(\psi(x)) \circ \langle x \circ !_B, \kappa_{x \in A}(\chi(x)) \circ \langle x \circ !_B, 1_B\rangle \rangle \\
                                                                                              &\underset{x}{=} \kappa_{x \in A}(\psi(x)) \circ \langle x \circ !_B, \chi(x)\rangle \\
                                                                                              &\underset{x}{=} \kappa_{x \in A}(\psi(x)) \circ \langle x \circ !_D, 1_D\rangle \circ \chi(x) &&(\textnormal{by $\spadesuit$})\\
                                                                                              &\underset{x}{=} \psi(x) \circ \chi(x),
            \end{align*}
            where $\spadesuit$ is the identity $\langle x \circ !_B, h \rangle = \langle x \circ !_D, 1_D\rangle \circ h$, for every $h : B \to D$, and the third and last equalities are by the inductive hypothesis.
            \item if $\varphi(x) = \langle \psi(x), \chi(x)\rangle$, with $\psi(x) : B \to D$, $\chi(x) : B \to E$ and $C = D \times E$, then:
            \begin{align*}
                \kappa_{x \in A}\langle \psi(x), \chi(x)\rangle \circ \langle x \circ !_B, 1_B\rangle &\underset{x}{=} \langle \kappa_{x \in A}\psi(x), \kappa_{x \in A}\chi(x)\rangle \circ \langle x \circ !_B, 1_B\rangle \\
                                                                                                      &\underset{x}{=} \langle \kappa_{x \in A}\psi(x) \circ \langle x \circ !_B, 1_B\rangle, \kappa_{x \in A}\chi(x) \circ \langle x \circ !_B, 1_B\rangle\rangle \\
                                                                                                      &\underset{x}{=} \langle \psi(x),\chi(x)\rangle;
            \end{align*}
            \item if $\varphi(x) = \psi(x)^*$, with $\psi(x) : B \times E \to D$ and $C = D^E$, then:
            \begin{align*}
                \kappa_{x \in A}(\psi(x)^*) \circ \langle x \circ !_B, 1_B\rangle &\underset{x}{=} (\kappa_{x \in A}\psi(x) \circ \alpha_{A,B,E})^* \circ \langle x \circ !_B, 1_B\rangle \\
                                                                                  &\underset{x}{=} (\kappa_{x \in A}\psi(x) \circ \alpha_{A,B,E} \circ (\langle x \circ !_B, 1_B\rangle \times 1_E))^* &&(\textnormal{by $\clubsuit$})\\
                                                                                  &\underset{x}{=} (\kappa_{x \in A}\psi(x) \circ \langle x \circ !_{B \times E}, 1_{B \times E}\rangle)^* &&(\textnormal{by $\heartsuit$})\\
                                                                                  &\underset{x}{=} \psi(x)^*,
            \end{align*}
            where $\clubsuit$ is the identity $f^* \circ g = (f \circ (g \times 1_E))^*$, for every $f : Y \times E \to D$ and $g : X \to Y$; $\heartsuit$ is the identity
            \[ \alpha_{A,B,E} \circ (\langle x \circ !_B, 1_B\rangle \times 1_E) = \langle x \circ !_{B \times E}, 1_{B \times E}\rangle, \]
            that can be proved using the definition of $\alpha_{A,B,E}$, and the last equality is by the inductive hypothesis.
        \end{itemize}
        \item[$\boxed{\textnormal{uniqueness}}$] Let $f : A \times B \to C$ be an arrow of $\mathcal C$ such that $f \circ \langle x \circ !_B, 1_B\rangle \underset{x}{=} \varphi(x)$, then:
        \begin{align*}
            \kappa_{x \in A}(\varphi(x)) &= \kappa_{x \in A}(f \circ \langle x \circ !_B, 1_B\rangle) \\
                                         &= \kappa_{x \in A}(f) \circ \langle \pi_{A,B}, \kappa_{x \in A}(\langle x \circ !_B, 1_B\rangle)\rangle \\
                                         &= (f \circ \pi_{A,A \times B}') \circ \langle \pi_{A,B}, \langle \pi_{A,1} \circ \langle \pi_{A,B}, !_B \circ \pi_{A,B}'\rangle, \pi_{A,B}' \rangle \rangle \\
                                         &= (f \circ \pi_{A,A \times B}') \circ \langle \pi_{A,B}, \langle \pi_{A,B}, \pi_{A,B}' \rangle \rangle \\
                                         &=f \circ (\pi_{A,A \times B}' \circ \langle \pi_{A,B}, \langle \pi_{A,B}, \pi_{A,B}' \rangle \rangle) \\
                                         &= f \circ \langle \pi_{A,B}, \pi_{A,B}' \rangle \underset{x}{=} f,
        \end{align*}
        where the last equality follows from $\langle \pi_{A,B}, \pi_{A,B}' \rangle \underset{x}{=} \langle \pi_{A,B} \circ 1_{A \times B}, \pi_{A,B}' \circ 1_{A \times B} \rangle \underset{x}{=} 1_{A \times B}$.
    \end{itemize}
\end{proof}

\begin{corollary}[Lambek's normal form]
    Given a cartesian closed category $\mathcal C$ and an indeterminate arrow $x : 1 \to A$, then for every polynomial arrow $\varphi(x) : 1 \to C$ there is a unique arrow $f : A \to C$ in $\mathcal C$ such that:
    \[ f \circ x \underset{x}{=} \varphi(x). \]
    Moreover, there's a unique $h : 1 \to C^A$ such that $\textnormal{ev}_{C,A} \circ \langle h,x \rangle \underset{x}{=} \varphi(x)$.
\end{corollary}

The idea behind the corollary is that, at the end of the day, in every cartesian closed category, both $\mathcal C$ and $\mathcal C[x]$, every arrow can be identified with an arrow from $1$.
So, thinking every arrow as a arrow from $1$, the functional completeness theorem above allows us to think every arrow of $\mathcal C[x]$ as an arrow of $\mathcal C$ composed with the indeterminate arrow.

\begin{proof}
    By functional completeness, there is a unique arrow $\kappa_{x \in A} \varphi(x) : A \times 1 \to C$ such that
    \[ \kappa_{x \in A} \varphi(x) \circ \langle x \circ !_1, 1_1\rangle \underset{x}{=} \varphi(x), \]
    since $\langle x \circ !_1, 1_1\rangle \underset{x}{=} \langle 1_A, !_A \rangle \circ x$, we can define $f := \kappa_{x \in A} \varphi(x) \circ \langle 1_A, !_A \rangle : A \to C$, hence:
    \[ f \circ x \underset{x}{=} \varphi(x). \]
    To prove uniqueness, let $g : A \to C$ be another arrow such that $g \circ x \underset{x}{=} \varphi(x)$, we have:
    \begin{align*}
        f &= \kappa_{x \in A}\varphi(x) \circ \langle 1_A,!_A\rangle \\
          &= \kappa_{x \in A}(g \circ x) \circ \langle 1_A,!_A\rangle \\
          &= \kappa_{x \in A}(g) \circ \langle \pi_{A,1},\kappa_{x \in A}(x)\rangle \circ \langle 1_A,!_A\rangle \\
          &= (g \circ \pi_{A,A}') \circ \langle \pi_{A,1},\pi_{A,1}\rangle \circ \langle 1_A,!_A\rangle \\
          &= g \circ (\pi_{A,A}' \circ \langle \pi_{A,1},\pi_{A,1}\rangle) \circ \langle 1_A,!_A\rangle \\
          &= g \circ \pi_{A,1} \circ \langle 1_A,!_A\rangle \\
          &= g.
    \end{align*}
    For the second part, we can define $h := \ulcorner f \urcorner \in \Hom(1,C^A)$, we saw in \Cref{prop:morphisms-as-global-elements} that $\ulcorner f \urcorner^\sharp = f$, so we have:
    \begin{align*}
        f \circ x &\underset{x}{=} h^\sharp \circ x \\    
                  &\underset{x}{=} \ev_{C,A} \circ \langle h \circ !_A,1_A \rangle \circ x \\
                  &\underset{x}{=} \ev_{C,A} \circ \langle h \circ 1_1, x \rangle &&(\textnormal{by the distributive law})\\
                  &\underset{x}{=} \ev_{C,A} \circ \langle h, x \rangle.
    \end{align*}
    The uniqueness of $h$ follows from the uniqueness of $f$ and the fact that $\ulcorner \cdot \urcorner$ is a bijection.
\end{proof}

\begin{remark}
    Once we introduce the $\lambda$-calculus, we will denote by $\lambda_{x \in A} \varphi(x)$ the unique arrow $h$ such that $h^\sharp x \underset{x}{=} \varphi(x)$.
\end{remark}
